\section{Conclusion and Future Work}
\label{sec:conclusion}

This work presented a research protocol and a reproducible analysis pipeline for an observational study on short-form video consumption, procrastination, and academic deadline adherence among university students~\cite{akniet2026shortform}. Two contributions were delivered. First, the protocol specifies a mixed-methods design---a 14-day diary, a baseline questionnaire, and device screen-time evidence---together with a pre-registered variable matrix and primary and guardrail metrics. Second, the pipeline was benchmarked on synthetic data: across 10--10{,}000 rows, a one-pass participant-level aggregation reduced median execution time by a factor of 1.85$\times$ to 4.26$\times$ relative to a grouped-aggregation-plus-merge baseline, with a 0\% error rate over 800 measured runs.

With respect to \textbf{RQ1} (descriptive), the protocol operationalizes the two target quantities---mean daily viewing time $X$ and baseline procrastination score $P$---and specifies their collection sources and validity checks. No empirical values are reported here: at the protocol stage, the repository contains only synthetic sample data. The descriptive question is therefore addressed at the level of measurement design, not of observed outcomes.

With respect to \textbf{RQ2 and RQ3} (associational), the protocol specifies the analytical machinery needed to test the core hypothesis \textbf{H1} (Spearman $\rho_s > 0$ between $X$ and $P$) and to estimate an adjusted incidence-rate ratio (IRR) for late submissions as a function of task-time and bedtime viewing. The benchmark confirms that the aggregation step underlying these statistics is fast and deterministic, but it does not test H1: the benchmark evaluates software performance, not student behaviour. Within the scope of this work, the results therefore \emph{suggest} readiness of the pipeline, not the presence or absence of the hypothesized associations.

Three concrete directions follow. \textbf{(i) Empirical data collection.} Execute the 14-day diary protocol with real participants and report $X$, $P$, the late-submission rate $R$, $\rho_s$, and the adjusted IRR against the pre-specified baseline (the ``neither'' viewing group). \textbf{(ii) Benchmark isolation.} Separate row count from participant-group count in the synthetic workload, and repeat timing on multiple machines and sessions, so that the observed speedup can be attributed to a specific factor rather than to their joint variation. \textbf{(iii) Pipeline extension.} Add end-to-end timing (CSV I/O and file writing) and an SQL-based aggregation baseline, to test whether the one-pass advantage persists outside the pandas implementation stack.

In summary, this repository delivers a reviewer-ready protocol and a validated aggregation pipeline. It does not claim behavioral findings. The next step is empirical data collection under the pre-registered design, followed by hypothesis testing for RQ2 and RQ3.
